\documentclass[12pt]{article}
\usepackage[utf8]{inputenc}
\usepackage[T1]{fontenc}
\usepackage{hyperref}
\title{Lycian A Language Resources for K-12 Curriculum Development}
\author{Compiled by Claude Sonnet 5.5 (claude-sonnet-5-5) for Andr\'as Kornai}
\date{}
\begin{document}
\maketitle

\section{Introduction}
Lycian A (ISO 639-3: xlc; Glottocode: lyci1241) is an extinct Anatolian Indo-European language of ancient Lycia in southwestern Anatolia. It survives in nearly 200 inscriptions, mostly on tombs and coins. Lycian B (Milyan) is a related variety. The Letoon trilingual inscription (Aramaic, Greek and Lycian), found in 1973, is a key text. The longer texts are still only partly understood.

\section{Primary Educational Resources}
I found no school curricula, primers or textbooks for Lycian. A heritage curriculum would have to be built from scholarly sources.

\section{Linguistic and Reference Materials}
\begin{itemize}
\item H. Craig Melchert, \emph{A Dictionary of the Lycian Language} (Beech Stave Press, 2004). It includes a grammatical sketch. \url{https://isdistribution.com/?p=1436030}
\item The publisher's PDF information sheet for the dictionary. \url{https://www.isdistribution.com/PDF/c6eae5d4-36c6-47a7-9eb3-af660fa429cb.pdf}
\item eDiAna, the Munich online dictionary of the minor languages of ancient Anatolia. Sample Lycian entry: \url{https://www.ediana-remake.gwi.uni-muenchen.de/dictionary/3014}
\item eDiAna project description. \url{https://wiki.digitalclassicist.org/EDiAna_–_Dictionary_of_the_Minor_Languages_of_Ancient_Anatolia}
\item Melchert, paper from Prague, 2019, on the interpretation of longer Lycian texts. \url{https://linguistics.ucla.edu/people/Melchert/melchertprague2019.pdf}
\item Lycian language (Wikipedia, ISO 639 xlc). \url{https://wikipedia.com/wiki/ISO_639:xlc}
\item Britannica, Anatolian languages using an alphabet. \url{https://britannica.com/topic/Anatolian-languages/Languages-using-an-alphabet}
\end{itemize}

\section{Religious and Community Resources}
No living community uses Lycian. The inscriptions are tomb and votive texts, including the Letoon trilingual.

\section{Digital and Audio Resources}
eDiAna offers online lexical entries that cite inscriptions. I found no audio resources.

\section{Curriculum Development Context}
Lycian is a classical-language case. It suits a heritage or ancient-languages unit that covers the script, the inscriptions and Lycia's history. The sources above are written for specialists. School materials would need to be created and adapted, and I found none that already exist.

\end{document}